\SourceLeaf{182}
\begin{equation*}
\frac{\gamma}{\alpha}=\frac{\pm\surd\bigl(D-\frac{a'A}{\alpha\alpha}\bigr)+b}{a'}
\quad \mathinner{\ldotp\ldotp\ldotp\ldotp\ldotp\ldotp\ldotp\ldotp} \tag*{[5]}
\end{equation*}
\begin{equation*}
\frac{\delta}{\beta}=\frac{\pm\surd\bigl(D+\frac{a'A'}{\beta\beta}\bigr)+b}{a'}
\quad \mathinner{\ldotp\ldotp\ldotp\ldotp\ldotp\ldotp\ldotp\ldotp} \tag*{[6]}
\end{equation*}
Aequatio 3, 4. 5, 6 reiicienda erit, si $\gamma,\delta,\alpha,\beta$ resp. $=0$. --- Sed dubium hic manet, quae \textit{signa} quantitatibus radicalibus tribui debeant; hoc sequenti modo decidemus.

Statim patet in [3] et [4] necessario signa superiora accipi debere, quando neque $\frac{\alpha}{\gamma}$ neque $\frac{\beta}{\delta}$ signum habeat signo ipsius $a$ oppositum; quoniam accepto signo inferiori $\frac{a\alpha}{\gamma}$ et $\frac{a\beta}{\delta}$ fierent quantitates negativae. Quia vero $A$ et $A'$ signa eadem habent, $\surd D$ cadet inter $\surd\bigl(D+\frac{aA}{\gamma\gamma}\bigr)$ et $\surd\bigl(D-\frac{aA'}{\delta\delta}\bigr)$ adeoque in hocce casu $\frac{\surd D-b}{a}$ inter $\frac{\alpha}{\gamma}$ et $\frac{\beta}{\delta}$. Quare pars prima theorematis pro casu priori est demonstrata.

Eodem modo perspicitur, in [5] et [6] necessario signa inferiora accipi debere, quando neque $\frac{\gamma}{\alpha}$ neque $\frac{\delta}{\beta}$ signum idem habeat ut $a'$ sive $a$, quia accepto superiori $\frac{a'\gamma}{\alpha}$, $\frac{a'\delta}{\beta}$ necessario fierent quantitates positivae. Unde protinus sequitur, $\frac{-\surd D+b}{a'}$ pro hocce casu iacere inter $\frac{\gamma}{\alpha}$ et $\frac{\delta}{\beta}$. Demonstrata est itaque etiam pars secunda theorematis pro casu posteriori. Quodsi aeque facile ostendi posset, in [3] et [4] signa inferiora accipi debere, quando neutra quantitatum $\frac{\alpha}{\gamma}$, $\frac{\beta}{\delta}$ signum idem habeat ut $a$, et in [5] et [6] superiora, quando neque $\frac{\gamma}{\alpha}$ neque $\frac{\delta}{\beta}$ signum oppositum habeat: hinc simili modo sequeretur, pro illo casu $\frac{-\surd D-b}{a}$ iacere inter $\frac{\alpha}{\gamma}$ et $\frac{\beta}{\delta}$, pro hoc $\frac{\surd D+b}{a'}$ inter $\frac{\gamma}{\alpha}$ et $\frac{\delta}{\beta}$, sive pars prima theorematis etiam pro casu posteriori, et secunda pro casu priori demonstratae forent. Sed quum illud difficile quidem non sit, attamen sine quibusdam ambagibus fieri nequeat. methodum sequentem praeferimus.

Quando nullus numerorum $\alpha,\beta.\gamma,\delta=0$, $\frac{\alpha}{\gamma}$ et $\frac{\beta}{\delta}$ eadem signa habebunt ut $\frac{\gamma}{\alpha}$, $\frac{\delta}{\beta}$. Quando itaque neutra harum quantitatum signum idem habet ut $a'$ sive $a$, adeoque $\frac{-\surd D+b}{a'}$ inter $\frac{\gamma}{\alpha}$ et $\frac{\delta}{\beta}$ cadit: neutra quantitatum $\frac{\alpha}{\gamma}$, $\frac{\beta}{\delta}$ signum idem ut $a$ habebit, cadetque $\frac{a'}{-\surd D+b}=\frac{-\surd D-b}{a}$ (propter $aa'=D-bb$) inter $\frac{\alpha}{\gamma}$ et $\frac{\beta}{\delta}$. Quare pro eo casu ubi neque $\alpha$ neque $\beta=0$, pars prima theor. etiam pro casu secundo est demonstrata (nam conditio ut neque $\gamma$ neque $\delta=0$, iam in theor. ipso est adiecta). Simili modo, quando nullus numerorum $\alpha,\beta,\gamma,\delta=0$,
\SourceLeaf{183}
et neque $\frac{\alpha}{\gamma}$ neque $\frac{\beta}{\delta}$ signum signo ipsius $a$ vel $a'$ oppositum habet, adeoque ${\frac{\surd D-b}{a'}}$ inter $\frac{\alpha}{\gamma}$ et $\frac{\beta}{\delta}$ iacet: etiam $\frac{\gamma}{\alpha}$ et $\frac{\delta}{\beta}$ signum oppositum signo ipsius $a'$ non habebit, cadetque $\frac{a}{\surd D-b}=\frac{\surd D+b}{a'}$ inter $\frac{\gamma}{\alpha}$ et $\frac{\delta}{\beta}$. In eo igitur casu ubi neque $\gamma$ neque $\delta=0$ pars secunda theor. etiam pro casu {secundo} est demonstrata.

Nihil itaque superesset quam ut demonstretur, partem primam theor. etiam pro casu secundo locum habere si alteruter numerorum $\alpha,\beta$ sit $=0$, et partem secundam pro casu primo si aut $\gamma$ aut $\delta=0$. At \textit{omnes hi casus sunt impossibiles}. Supponamus enim, pro parte \textit{prima} theor., esse neque $\gamma$ neque $\delta=0$; $\frac{\alpha}{\gamma}$, $\frac{\beta}{\delta}$ non habere signum idem ut $a$ atque esse 1) $\alpha=0$. Tum ex aequ. $\alpha\delta-\beta\gamma=\pm1$ fit $\beta=\pm1$, $\gamma=\pm1$. Hinc ex [1] $A=-a'$, quare $A$ et $a'$, adeoque etiam $a$ et $A'$ signa opposita habent, unde fit $\surd\bigl(D-\frac{aA'}{\delta\delta}\bigr)>\surd D>b$. Hinc patet in [4] necessario signum inferius accipi debere, quia accepto superiori $\frac{\beta}{\delta}$ manifesto signum idem obtineret ut $a$. Fit itaque $\frac{\beta}{\delta}>\frac{-\surd D-b}{a}>1$ (propter $a<\surd D+b$ ex def. formae reductae), \textit{Q. E. A.} quum $\beta=\pm1$, et $\delta$ non $=0$. --- 2) Sit $\beta=0$. Tum ex aequ. $\alpha\delta-\beta\gamma=\pm1$ fit $\alpha=\pm1$, $\delta=\pm1$. Hinc ex [2] $-A'=-a'$, quare $a'$ et $a$ et $A$ signa eadem habebunt, unde fit ${\surd\bigl(D+\frac{aA}{\alpha\alpha}\bigr)}>\surd D>b$. Hinc patet in [3] signum inferius accipi debere, quia accepto superiori $\frac{\alpha}{\gamma}$ signum idem obtineret ut $a$. Fit itaque $\frac{\alpha}{\gamma}>\frac{-\surd D-b}{a}>1$, \textit{Q. E. A.} eadem ratione ut ante. --- Pro parte \textit{secunda} si supponimus neque $\alpha$ neque $\beta=0$; $\frac{\gamma}{\alpha}$, $\frac{\delta}{\beta}$ non habere signum signo ipsius $a'$ oppositum atque 1) $\gamma=0$: ex aequ. $\alpha\delta-\beta\gamma=\pm1$ fit $\alpha=\pm1$, $\delta=\pm1$. Hinc ex [1] $A=a$, quare $a'$ et $A'$ signa eadem habebunt, unde fit $\surd\bigl(D+\frac{a'A'}{\beta\beta}\bigr)>\surd D>b$. Quocirca in [6] signum superius erit accipiendum, quia accepto inferiori, $\frac{\delta}{\beta}$ obtineret signum oppositum signo ipsius $a'$. Fit igitur $\frac{\delta}{\beta}>\frac{\surd D+b}{a'}>1$, \textit{Q. E. A.}, quia $\delta=\pm1$ et $\beta$ non $=0$. Tandem 2) si esset $\delta=0$, ex $\alpha\delta-\beta\gamma=\pm1$ fit $\beta=\pm1$, $\gamma=\pm1$, adeoque ex [2] $-A'=a$. Hinc $\surd\bigl(D-\frac{a'A}{\alpha\alpha}\bigr)>\surd D>b$, quare in [5] signum superius accipiendum. Hinc $\frac{\gamma}{\alpha}>\frac{\surd D+b}{a'}>1$, \textit{Q. E. A.} --- Quare theorema in omni sua extensione est demonstratum.

Quum differentia inter $\frac{\alpha}{\gamma}$ et $\frac{\beta}{\delta}$ sit $=\frac{1}{\gamma\delta}$: differentia inter $\frac{\pm\surd D-b}{a}$ et $\frac{\alpha}{\gamma}$ vel $\frac{\beta}{\delta}$ erit $<\frac{1}{\gamma\delta}$; inter $\frac{\pm\surd D-b}{a}$ autem et $\frac{\alpha}{\gamma}$, vel inter illam quantitatem et $\frac{\beta}{\delta}$ nulla fractio iacere poterit, cuius denominator non sit maior quam $\gamma$ aut $\delta$ (\textit{lemma praec.}). --- Eodem modo differentia quantitatis ${\frac{\pm\surd D+b}{a}}$ a fractione $\frac{\gamma}{\alpha}$ vel hac $\frac{\delta}{\beta}$
\SourceLeaf{184}
erit minor quam $\frac{1}{\alpha\beta}$, et inter illam quantitatem et neutram harum fractionum iacere potest fractio cuius denominator non sit maior quam $\alpha$ et $\beta$.

\Needspace{8\baselineskip}
\Article{192}
Ex applicatione theor. praec. ad algorithmum art. 188 sequitur, quantitatem $\frac{\surd D-b}{a}$ quam per $L$ designabimus, iacere inter $\frac{\alpha'}{\gamma'}$ et $\frac{\beta'}{\delta'}$; inter $\frac{\alpha''}{\gamma''}$ et $\frac{\beta''}{\delta''}$; inter $\frac{\alpha'''}{\gamma'''}$ et $\frac{\beta'''}{\delta'''}$ etc. (facile enim ex art. 189, 3 fin. deducitur, nullum horum limitum habere signum oppositum signo ipsius $a$; quare quantitati radicali $\surd D$ signum positivum tribui debet) sive inter $\frac{\alpha'}{\gamma'}$ et $\frac{\alpha''}{\gamma''}$; inter $\frac{\alpha''}{\gamma''}$ et $\frac{\alpha'''}{\gamma'''}$ etc. Omnes itaque fractiones $\frac{\alpha'}{\gamma'}$, $\frac{\alpha'''}{\gamma'''}$, $\frac{\alpha'''''}{\gamma'''''}$ etc. ipsi $L$ ab eadem parte iacebunt, omnesque $\frac{\alpha''}{\gamma''}$, $\frac{\alpha''''}{\gamma''''}$, $\frac{\alpha''''''}{\gamma''''''}$ etc. a parte altera. Quoniam vero $\gamma'<\gamma'''$, $\frac{\alpha'}{\gamma'}$ iacebit \textit{extra} $\frac{\alpha'''}{\gamma'''}$ et $L$, similique ratione $\frac{\alpha''}{\gamma''}$ extra $L$ et $\frac{\alpha''''}{\gamma''''}$; $\frac{\alpha'''}{\gamma'''}$ extra $L$ et $\frac{\alpha'''''}{\gamma'''''}$ etc. Unde manifestum est, has quantitates iacere sequenti ordine:
\[
\frac{\alpha'}{\gamma'},\ \frac{\alpha'''}{\gamma'''},\ \frac{\alpha'''''}{\gamma'''''}\mathinner{\ldotp\ldotp\ldotp} L\mathinner{\ldotp\ldotp\ldotp}\frac{\alpha''''''}{\gamma''''''},\ \frac{\alpha''''}{\gamma''''},\ \frac{\alpha''}{\gamma''}
\]
Differentia autem inter $\frac{\alpha'}{\gamma'}$ et $L$ erit minor quam differentia inter $\frac{\alpha'}{\gamma'}$ et $\frac{\alpha''}{\gamma''}$ \textit{i. e.} $<\frac{1}{\gamma'\gamma''}$, similique ratione differentia inter $\frac{\alpha''}{\gamma''}$ et $L$ erit $<\frac{1}{\gamma''\gamma'''}$ etc. Quamobrem fractiones $\frac{\alpha'}{\gamma'}$, $\frac{\alpha''}{\gamma''}$, $\frac{\alpha'''}{\gamma'''}$ etc. continuo propius ad limitem $L$ accedunt, et quoniam $\gamma',\gamma'',\gamma'''$ continuo in infinitum crescunt, differentia fractionum a limite quavis quantitate data minor fieri potest.

Ex art. 189 nulla quantitatum $\frac{\gamma}{\alpha}$, $\frac{{}'\gamma}{{}'\alpha}$, $\frac{{}''\gamma}{{}''\alpha}$ etc. signum idem habebit ut $a$; hinc per ratiocinia praecedentibus omnino similia sequitur, illas et hanc $\frac{-\surd D+b}{a'}$, quam per $L'$ designabimus, iacere sequenti ordine:
\[
\frac{\gamma}{\alpha},\ \frac{{}''\gamma}{{}''\alpha},\ \frac{{}''''\gamma}{{}''''\alpha}\mathinner{\ldotp\ldotp\ldotp\ldotp\ldotp} L'\mathinner{\ldotp\ldotp\ldotp\ldotp\ldotp}\frac{{}'''''\gamma}{{}'''''\alpha},\ \frac{{}'''\gamma}{{}'''\alpha},\ \frac{{}'\gamma}{{}'\alpha}
\]
Differentia autem inter $\frac{\gamma}{\alpha}$ et $L'$ minor erit quam $\frac{1}{{}'\alpha\alpha}$, differentia inter $\frac{{}'\gamma}{{}'\alpha}$ et $L'$ minor quam $\frac{1}{{}''\alpha{}'\alpha}$ etc. Quare fractiones $\frac{\gamma}{\alpha}$, $\frac{{}'\gamma}{{}'\alpha}$ etc. continuo propius ad $L'$ accedent, et differentia quavis quantitate data minor fieri poterit.

In \textit{ex.} art. 188 fit $L=\frac{\surd79-{8}}{3}=0{,}2960648$ et fractiones appropinquantes $\frac{0}{1}$, $\frac{1}{3}$, $\frac{2}{7}$, $\frac{3}{10}$, $\frac{5}{17}$, $\frac{8}{27}$, $\frac{45}{152}$, $\frac{143}{483}$, etc. Est autem $\frac{143}{483}=0{,}2960662$. --- Ibidem fit $L'=\frac{-\surd79+{8}}{5}=-0{,}1776388$, fractionesque approximantes $\frac{0}{1}$, $-\frac{1}{5}$, $-\frac{1}{6}$, $-\frac{2}{11}$, $-\frac{3}{17}$, $-\frac{8}{45}$, $-\frac{27}{152}$, $-\frac{143}{805}$ etc. Est vero $\frac{143}{805}=0{,}1776397$.
\SourceLeaf{185}
\Needspace{8\baselineskip}
\Article{193}
\textsc{Theorema.} \textit{Si formae reductae $f$, $F$ proprie aequivalentes sunt: altera in alterius periodo contenta erit.}

Sit $f=(a,b,-a')$, $F=(A,B,-A')$, determinans harum formarum $D$, transeatque illa in hanc per substitutionem propriam $\mathfrak{A},\mathfrak{B},\mathfrak{C},\mathfrak{D}$. Tum dico, si periodus formae $f$ quaeratur progressioque utrinque infinita formarum reductarum atque transformationum formae $f$ in ipsas eruatur, eodem modo ut art. 188: \textit{vel} $+\mathfrak{A}$ fore aequalem termino alicui progressionis $\mathinner{\ldotp\ldotp\ldotp}{}''\alpha,{}'\alpha,\alpha,\alpha',\alpha''\mathinner{\ldotp\ldotp\ldotp}$, hocque posito $=\alpha^m$, $+\mathfrak{B}$ fore $=\beta^m$, $+\mathfrak{C}=\gamma^m$, $+\mathfrak{D}=\delta^m$; \textit{vel} $-\mathfrak{A}$ fore aequalem termino alicui $\alpha^m$, et $-\mathfrak{B},-\mathfrak{C},-\mathfrak{D}$ resp. $=\beta^m,\gamma^m,\delta^m$ (ubi $m$ etiam indicem negativum designare potest). In utroque casu $F$ manifesto identica erit cum $f^m$.

\textit{Dem.} I. Habentur quatuor aequationes,
\begin{align*}
a\mathfrak{A}\mathfrak{A}+2b\mathfrak{A}\mathfrak{C}-a'\mathfrak{C}\mathfrak{C}&=A \quad \mathinner{\ldotp\ldotp\ldotp\ldotp\ldotp\ldotp\ldotp} \tag*{[1]}\\
a\mathfrak{A}\mathfrak{B}+b(\mathfrak{A}\mathfrak{D}+\mathfrak{B}\mathfrak{C})-a'\mathfrak{C}\mathfrak{D}&=B \quad \mathinner{\ldotp\ldotp\ldotp\ldotp\ldotp\ldotp} \tag*{[2]}\\
a\mathfrak{B}\mathfrak{B}+2b\mathfrak{B}\mathfrak{D}-a'\mathfrak{D}\mathfrak{D}&=-A' \quad \mathinner{\ldotp\ldotp\ldotp\ldotp\ldotp\ldotp} \tag*{[3]}\\
\mathfrak{A}\mathfrak{D}-\mathfrak{B}\mathfrak{C}&=1 \quad \mathinner{\ldotp\ldotp\ldotp\ldotp\ldotp\ldotp\ldotp\ldotp\ldotp\ldotp} \tag*{[4]}
\end{align*}
Consideramus autem \textit{primo} casum, ubi aliquis numerorum $\mathfrak{A}\ \mathfrak{B},\mathfrak{C},\mathfrak{D}=0$.

$1^\circ$ Si $\mathfrak{A}=0$, fit ex [4] $\mathfrak{B}\mathfrak{C}=-1$, adeoque $\mathfrak{B}=\pm1$, $\mathfrak{C}=\mp1$. Hinc ex [1], $-a'=A$; ex [2], $-b\pm a'\mathfrak{D}=B$ sive $B\equiv-b\ (\mathrm{mod.}\ a'\text{ vel }A)$; unde sequitur formam $(A,B,-A')$ formae $(a,b,-a')$ ab ultima parte contiguam esse. Quoniam vero illa est reducta, necessario cum $f'$ identica erit. Ergo $B=b'$, adeoque ex [2] $b+b'=-a'\mathfrak{C}\mathfrak{D}=\pm a'\mathfrak{D}$; hinc propter $\frac{b+b'}{-a'}=h'$, fit $\mathfrak{D}=\mp h'$. Unde colligitur, $\mp\mathfrak{A},\mp\mathfrak{B},\mp\mathfrak{C},\mp\mathfrak{D}$ esse resp. $=0,-1,+1,h'$ sive $=\alpha',\beta',\gamma',\delta'$.

$2^\circ$ Si $\mathfrak{B}=0$, fit ex [4] $\mathfrak{A}=\pm1$, $\mathfrak{D}=\pm1$; ex [3] $a'=A'$; ex [2] $b\mp a'\mathfrak{C}=B$, sive $b\equiv B\ (\mathrm{mod.}\ a')$. Quoniam vero tum $f$ tum $F$ sunt formae reductae: tum $b$ tum $B$ iacebunt inter $\surd D$ et $\surd D\mp a'$ (prout $a'$ pos. vel neg., art. 185, 5). Quare erit necessario $b=B$, et $\mathfrak{C}=0$. Hinc formae $f$, $F$ sunt identicae atque $\pm\mathfrak{A},\pm\mathfrak{B},\pm\mathfrak{C},\pm\mathfrak{D}=1,0,0,1=\alpha,\beta,\gamma\ \delta$ (resp.).

$3^\circ$ Si $\mathfrak{C}=0$, fit ex [4] $\mathfrak{A}=\pm1$, $\mathfrak{D}=\pm1$; ex [1] $a=A$; ex [2] $\pm a\mathfrak{B}+b=B$ sive $b\equiv B\ (\mathrm{mod.}\ a)$. Quia vero tum $b$ tum $B$ iacent inter $\surd D$
\SourceLeaf{186}
et $\surd D\mp a$: erit necessario $B=b$ et $\mathfrak{B}=0$. Quare casus hic a praecedente non differt.

$4^\circ$ Si $\mathfrak{D}=0$, fit ex [4] $\mathfrak{B}=\pm1$, $\mathfrak{C}=\mp1$; ex [3] $a=-A'$; ex [2] $\pm a\mathfrak{A}-b=B$ sive $B\equiv-b\ (\mathrm{mod.}\ a)$. Hinc forma $F$ formae $f$ a parte prima contigua erit, et proin cum forma ${}'f$ identica. Quare propter $\frac{{}'b+b}{a}=h$, et $B={}'b$, erit $\pm\mathfrak{A}=h$. Unde colligitur $\pm\mathfrak{A},\pm\mathfrak{B},\pm\mathfrak{C},\pm\mathfrak{D}$ resp. esse $=h,1,-1,0$, $={}'\alpha,{}'\beta,{}'\gamma,{}'\delta$.

Superest itaque casus ubi nullus numerorum $\mathfrak{A},\mathfrak{B},\mathfrak{C},\mathfrak{D}=0$. Hic per Lemma art. 190 quantitates $\frac{\mathfrak{A}}{\mathfrak{C}}$, $\frac{\mathfrak{B}}{\mathfrak{D}}$, $\frac{\mathfrak{C}}{\mathfrak{A}}$, $\frac{\mathfrak{D}}{\mathfrak{B}}$ idem signum habebunt, oriunturque inde duo casus, quum signum hoc vel cum signo ipsorum $a$, $a'$ convenire vel ipsi oppositum esse possit.

II. Si $\frac{\mathfrak{A}}{\mathfrak{C}}$, $\frac{\mathfrak{B}}{\mathfrak{D}}$ idem signum habent ut $a$: quantitas $\frac{\surd D-b}{a}$ (quam designabimus per $L$) inter has fractiones sita erit (art. 191). Demonstrabimus iam, $\frac{\mathfrak{A}}{\mathfrak{C}}$ aequalem fore alicui fractionum $\frac{\alpha''}{\gamma''}$, $\frac{\alpha'''}{\gamma'''}$, $\frac{\alpha''''}{\gamma''''}$ etc., atque $\frac{\mathfrak{B}}{\mathfrak{D}}$ proxime sequenti, scilicet si $\frac{\mathfrak{A}}{\mathfrak{C}}$ fuerit $=\frac{\alpha^m}{\gamma^m}$, $\frac{\mathfrak{B}}{\mathfrak{D}}$ fore $=\frac{\alpha^{m+1}}{\gamma^{m+1}}$. In art. praec. ostendimus, quantitates $\frac{\alpha'}{\gamma'}$, $\frac{\alpha''}{\gamma''}$, $\frac{\alpha'''}{\gamma'''}$ etc. (quas brevitatis gratia per (1), (2), (3) etc. denotabimus), atque $L$, hunc ordinem (I) observare: $(1),(3),(5)\mathinner{\ldotp\ldotp\ldotp} L\mathinner{\ldotp\ldotp\ldotp}(6),(4),(2)$; prima harum quantitatum est $=0$ (propter $\alpha'=0$), reliquae omnes idem signum habent ut $L$ sive $a$. Quoniam vero per hyp. $\frac{\mathfrak{A}}{\mathfrak{C}}$, $\frac{\mathfrak{B}}{\mathfrak{D}}$ (pro quibus scribemus $\mathfrak{M},\mathfrak{N}$) idem signum habent: patet has quantitates ipsi (1) a dextra iacere (aut si mavis ab eadem parte a qua $L$), et quidem, quum $L$ iaceat inter ipsas, alteram ipsi $L$ a dextra, alteram a laeva. Facile vero ostendi potest, $\mathfrak{M}$ ipsi (2) a dextra iacere non posse, alioquin enim $\mathfrak{N}$ iaceret inter (1) et $L$, unde sequeretur \textit{primo} (2) iacere inter $\mathfrak{M}$ et $\mathfrak{N}$, adeoque denominatorem fractionis (2) maiorem esse denominatore fractionis $\mathfrak{N}$ (art. 190), \textit{secundo} $\mathfrak{N}$ iacere inter (1) et (2), adeoque denom. fractionis $\mathfrak{N}$ esse maiorem quam denom. fractionis (2), \textit{Q. E. A.}

Supponamus $\mathfrak{M}$ nulli fractionum (2), (3), (4) etc. aequalem esse, ut, quid inde sequatur, videamus. Tum manifestum est, si fractio $\mathfrak{M}$ ipsi $L$ a laeva iaceat, necessario eam sitam esse aut inter (1) et (3), aut inter (3) et (5), aut inter (5) et (7) etc. (quoniam $L$ est irrationalis, adeoque ipsi $\mathfrak{M}$ certo inaequalis, fractionesque (1), (3), (5) etc. quavis quantitate data, ipsi $L$ inaequali, propius ad $L$ accedere possunt). Si vero $\mathfrak{M}$ ipsi $L$ a dextra iacet: necessario iacebit aut inter (2) et (4), aut inter (4) et (6) aut inter (6) et (8) etc. Ponamus itaque $\mathfrak{M}$ iacere inter
\SourceLeaf{187}
$(m)$ et $(m+2)$, patetque quantitates $\mathfrak{M},(m),(m+1),(m+2),L$ iacere sequenti ordine,
\begin{center}
(II)\HistoricalNote{*)}{Nihil hic refert, sive ordo in (II) idem sit ut in (I), sive huic oppositus, \textit{i. e.} sive $(m)$ etiam in (I) ipsi $L$ a laeva iaceat sive a dextra.}:\qquad $(m),\ (\mathfrak{M}),\ (m+2),\ L,\ (m+1)$
\end{center}
Tum erit necessario $\mathfrak{N}=(m+1)$. Iacebit enim $\mathfrak{N}$ ipsi $L$ a dextra; si vero etiam ipsi $(m+1)$ a dextra iaceret, $(m+1)$ iaceret inter $\mathfrak{M}$ et $\mathfrak{N}$, unde $\gamma^{m+1}>\mathfrak{C}$, $\mathfrak{M}$ vero inter $(m)$ et $(m+1)$ unde $\mathfrak{C}>\gamma^{m+1}$ (art. 190), \textit{Q. E. A.}; si vero $\mathfrak{N}$ ipsi $(m+1)$ a laeva iaceret, sive inter $(m+2)$ et $(m+1)$, foret $\mathfrak{D}>\gamma^{m+2}$, et quia $(m+2)$ inter $\mathfrak{M}$ et $\mathfrak{N}$, foret $\gamma^{m+2}>\mathfrak{D}$, \textit{Q. E. A.} Erit itaque $\mathfrak{N}=(m+1)$, sive $\frac{\mathfrak{B}}{\mathfrak{D}}=\frac{\alpha^{m+1}}{\gamma^{m+1}}=\frac{\beta^m}{\delta^m}$.

Quia $\mathfrak{A}\mathfrak{D}-\mathfrak{B}\mathfrak{C}=1$, $\mathfrak{B}$ erit primus ad $\mathfrak{D}$ et ex simili ratione $\beta^m$ primus ad $\delta^m$. Unde facile perspicitur aequationem $\frac{\mathfrak{B}}{\mathfrak{D}}=\frac{\beta^m}{\delta^m}$ consistere non posse, nisi fuerit aut $\mathfrak{B}=\beta^m$, $\mathfrak{D}=\delta^m$, aut $\mathfrak{B}=-\beta^m$, $\mathfrak{D}=-\delta^m$. Iam quum forma $f$ per substitutionem propriam $\alpha^m,\beta^m,\gamma^m,\delta^m$ in formam $f^m$ transmutetur, quae est $(\pm a^m,b^m,\mp a^{m+1})$: habebuntur aequationes
\begin{align*}
a\alpha^m\alpha^m+2b\alpha^m\gamma^m-a'\gamma^m\gamma^m&=\pm a^m \quad \mathinner{\ldotp\ldotp\ldotp\ldotp\ldotp} \tag*{[5]}\\
a\alpha^m\beta^m+b(\alpha^m\delta^m+\beta^m\gamma^m)-a'\gamma^m\delta^m&=b^m \quad \mathinner{\ldotp\ldotp\ldotp\ldotp} \tag*{[6]}\\
a\beta^m\beta^m+2b\beta^m\delta^m-a'\delta^m\delta^m&=\mp a^{m+1} \quad \mathinner{\ldotp\ldotp\ldotp\ldotp\ldotp} \tag*{[7]}\\
\alpha^m\delta^m-\beta^m\gamma^m&=1 \quad \mathinner{\ldotp\ldotp\ldotp\ldotp\ldotp\ldotp\ldotp\ldotp\ldotp} \tag*{[8]}
\end{align*}
Hinc fit: (ex aequ. 7 et 3), $\mp a^{m+1}=-A'$. Porro multiplicando aequationem [2] per $\alpha^m\delta^m-\beta^m\gamma^m$, aequationem [6] per $\mathfrak{A}\mathfrak{D}-\mathfrak{B}\mathfrak{C}$ et subtrahendo facile per evolutionem confirmatur esse
\begin{equation*}
\begin{aligned}
B-b^m={}&(\mathfrak{C}\alpha^m-\mathfrak{A}\gamma^m)\bigl(a\mathfrak{B}\beta^m\\
&\qquad+b(\mathfrak{D}\beta^m+\mathfrak{B}\delta^m)-a'\mathfrak{D}\delta^m\bigr)\\
&+(\mathfrak{B}\delta^m-\mathfrak{D}\beta^m)\bigl(a\mathfrak{A}\alpha^m\\
&\qquad+b(\mathfrak{C}\alpha^m+\mathfrak{A}\gamma^m)-a'\mathfrak{C}\gamma^m\bigr)
\quad \mathinner{\ldotp\ldotp}
\end{aligned}
\tag*{[9]}
\end{equation*}
sive quoniam vel $\beta^m=\mathfrak{B}$, $\delta^m=\mathfrak{D}$ vel $\beta^m=-\mathfrak{B}$, $\delta^m=-\mathfrak{D}$,
\[
\begin{aligned}
B-b^m&=\pm(\mathfrak{C}\alpha^m-\mathfrak{A}\gamma^m)(a\mathfrak{B}\mathfrak{B}+2b\mathfrak{B}\mathfrak{D}-a'\mathfrak{D}\mathfrak{D})\\
&=\mp(\mathfrak{C}\alpha^m-\mathfrak{A}\gamma^m)A'
\end{aligned}
\]
Hinc $B\equiv b^m\ (\mathrm{mod.}\ A')$; quia vero tum $B$ tum $b^m$, inter $\surd D$ et $\surd D\mp A'$ iacent, necessario erit $B=b^m$ adeoque $\mathfrak{C}\alpha^m-\mathfrak{A}\gamma^m=0$, sive $\frac{\mathfrak{A}}{\mathfrak{C}}=\frac{\alpha^m}{\gamma^m}$, \textit{i. e.} $\mathfrak{M}=(m)$.

Hoc modo itaque ex suppositione, $\mathfrak{M}$ nulli quantitatum (2), (3), (4) etc. aequalem esse, deduximus, eam revera alicui aequalem esse. Quodsi vero ab
\SourceLeaf{188}
initio supponimus, esse $\mathfrak{M}=(m)$, manifesto erit vel $\mathfrak{A}=\alpha^m$, $\mathfrak{C}=\gamma^m$, vel $-\mathfrak{A}=\alpha^m$, $-\mathfrak{C}=\gamma^m$. In utroque casu fit ex [1] et [5] $A=\pm a^m$, et ex [9] $B-b^m=\pm(\mathfrak{B}\delta^m-\mathfrak{D}\beta^m)A$, sive $B\equiv b^m\ (\mathrm{mod.}\ A)$. Hinc simili modo ut supra concluditur $B=b^m$, et hinc $\mathfrak{B}\delta^m=\mathfrak{D}\beta^m$; quare quum $\mathfrak{B}$ ad $\mathfrak{D}$ primus sit et $\beta^m$ ad $\delta^m$: erit aut $\mathfrak{B}=\beta^m$, $\mathfrak{D}=\delta^m$ aut $-\mathfrak{B}=\beta^m$, $-\mathfrak{D}=\delta^m$, et proin ex [7] $-A'=\mp a^{m+1}$. Quamobrem formae $F$, $f^m$ identicae erunt. Adiumento aequationis $\mathfrak{A}\mathfrak{D}-\mathfrak{B}\mathfrak{C}=\alpha^m\delta^m-\beta^m\gamma^m$ autem nullo negotio probatur, poni debere $+\mathfrak{B}=\beta^m$, $+\mathfrak{D}=\delta^m$, quando $+\mathfrak{A}=\alpha^m$, $+\mathfrak{C}=\gamma^m$; contra $-\mathfrak{B}=\beta^m$, ${-\mathfrak{D}=-\delta^m}$, quando $-\mathfrak{A}=\alpha^m$, $-\mathfrak{C}=\gamma^m$. \textit{Q. E. D.}

III. Si signum quantitatum $\frac{\mathfrak{A}}{\mathfrak{C}}$ etc. signo ipsius $a$ oppositum: demonstratio praecedenti tam similis est, ut praecipua tantum momenta addigitavisse sufficiat. Iacebit $\frac{-\surd D+b}{a'}$ inter $\frac{\mathfrak{C}}{\mathfrak{A}}$ et $\frac{\mathfrak{D}}{\mathfrak{B}}$. Fractio $\frac{\mathfrak{D}}{\mathfrak{B}}$ alicui fractionum
\begin{equation*}
\frac{{}'\delta}{{}'\beta},\ \frac{{}''\delta}{{}''\beta},\ \frac{{}'''\delta}{{}'''\beta}\text{ etc. aequalis erit}
\quad \mathinner{\ldotp\ldotp\ldotp\ldotp\ldotp\ldotp\ldotp} \tag*{(I)}
\end{equation*}
\begin{equation*}
\text{qua posita }=\frac{{}^m\delta}{{}^m\beta},\qquad \frac{\mathfrak{C}}{\mathfrak{A}}\text{ erit }=\frac{{}^m\gamma}{{}^m\alpha}
\quad \mathinner{\ldotp\ldotp\ldotp\ldotp\ldotp\ldotp\ldotp\ldotp\ldotp\ldotp\ldotp} \tag*{(II)}
\end{equation*}
Demonstratur autem (I) ita: Si $\frac{\mathfrak{D}}{\mathfrak{B}}$ nulli illarum fractionum aequalis esse supponitur: inter duas tales $\frac{{}^m\delta}{{}^m\beta}$ et $\frac{{}^{m+2}\delta}{{}^{m+2}\beta}$ iacere debebit. Hinc vero eodem modo ut supra deducitur, necessario esse
\[
\frac{\mathfrak{C}}{\mathfrak{A}}=\frac{{}^{m+1}\delta}{{}^{m+1}\beta}=\frac{{}^m\gamma}{{}^m\alpha}
\]
atque vel $\mathfrak{A}={}^m\alpha$, $\mathfrak{C}={}^m\gamma$, vel $-\mathfrak{A}={}^m\alpha$, $-\mathfrak{C}={}^m\gamma$. Quoniam vero $f$ per substitutionem propriam ${}^m\alpha,{}^m\beta,{}^m\gamma,{}^m\delta$ in formam
\[
{}^m f=(\pm{}^m a,{}^m b,\mp{}^{m-1}a)
\]
transit: hinc emergunt tres aequationes, ex quibus coniunctis cum aequ. 1, 2, 3, 4 atque hac, ${}^m\alpha{}^m\delta-{}^m\beta{}^m\gamma=1$ deducitur eodem modo ut supra, terminum primum $A$ formae $F$ termino primo formae ${}^m f$ aequalem esse, illiusque terminum medium medio huius congruum secundum modulum $A$, unde sequitur, quia utraque forma est reducta, adeoque utriusque terminus medius inter $\surd D$ et $\surd D\mp A$ situs, hos terminos medios aequales esse: hinc vero deducitur $\frac{{}^m\delta}{{}^m\beta}=\frac{\mathfrak{D}}{\mathfrak{B}}$. Veritas itaque assertionis (I) derivata hic est ex suppositione illam esse falsam.

Supponendo autem $\frac{{}^m\delta}{{}^m\beta}=\frac{\mathfrak{D}}{\mathfrak{B}}$, prorsus simili modo et per easdem aequatio%
\SourceLeaf{189}%
nes demonstratur, esse etiam ${\frac{{}^m\gamma}{{}^m\beta}}=\frac{\mathfrak{C}}{\mathfrak{A}}$, quod erat secundum (II). Hinc vero adiumento aequationum $\mathfrak{A}\mathfrak{D}-\mathfrak{B}\mathfrak{C}=1$, ${}^m\alpha{}^m\delta-{}^m\beta{}^m\gamma=1$ deducitur esse vel
\[
\mathfrak{A}={}^m\alpha,\ \mathfrak{B}={}^m\beta,\ \mathfrak{C}={}^m\gamma,\ \mathfrak{D}={}^m\delta
\]
vel
\[
-\mathfrak{A}={}^m\alpha,\ -\mathfrak{B}={}^m\beta,\ -\mathfrak{C}={}^m\gamma,\ -\mathfrak{D}={}^m\delta
\]
formasque $F$, ${}^m f$ identicas. \textit{Q. E. D.}

\Needspace{8\baselineskip}
\Article{194}
Quum formae quas supra socias vocavimus (art. 187, 6) semper sint improprie aequivalentes (art. 159), perspicuum est, si formae reductae $F$, $f$ improprie aequivalentes sint, formaeque $F$ socia forma $G$, formas $f$, $G$ proprie aequivalentes fore adeoque formam $G$ in periodo formae $f$ contentam. Quodsi itaque formae $F$, $f$ tum proprie tum improprie aequivalentes sunt, patet, tum $F$ tum $G$ in periodo formae $f$ reperiri debere. Quare periodus haec sibi ipsi socia erit, duasque formas ancipites continebit (art. 187, 7). Vnde theorema art. 165 egregie confirmatur, ex quo iam poteramus esse certi, formam aliquam ancipitem dari formis $F$, $f$ aequivalentem.

\Needspace{8\baselineskip}
\Article{195}
\textsc{Problema.} \textit{Propositis duabus formis quibuscunque $\Phi$, $\varphi$ eiusdem determinantis: diiudicare utrum aequivalentes sint, annon.}

\textit{Sol.} Quaerantur duae formae reductae $F$, $f$ propositis $\Phi$, $\varphi$ resp. proprie aequivalentes (art. 183). Quae prout aut proprie tantum aequivalent, aut improprie tantum, aut utroque modo, aut neutro; etiam propositae aut proprie tantum aequivalentes erunt, aut improprie tantum, aut utroque aut neutro modo. Evolvatur periodus alterutrius formae reductae \textit{e. g.} periodus formae $f$. Si forma $F$ in hac periodo occurrit neque vero simul forma ipsi $F$ socia, manifesto casus \textit{primus} locum habebit; contra si socia haec adest neque vero $F$ ipsa, \textit{secundus}; si utraque, \textit{tertius}; si neutra, \textit{quartus}.

\textit{Ex.} Propositae sint formae $(129,92,65)$, $(42,59,81)$ determinantis 79. His proprie aequivalentes inveniuntur reductae $(10,7,-3)$, $(5,8,-3)$. Periodus formae prioris haec est: $(10,7,-3)$, $(-3,8,5)$, $(5,7,-6)$, $(-6,5,9)$, $(9,4.\,-7)$, $(-7,3,10)$. In qua quum forma $(5,8,-3)$ ipsa non reperiatur,
\SourceLeaf{190}
sed tamen socia $(-3,8,5)$: formas propositas improprie tantum aequivalere concludimus.

Si omnes formae reductae determinantis dati eodem modo ut supra (art. 187, 5) in periodos $P$, $Q$, $R$ etc. distribuuntur, atque e quavis periodo forma aliqua ad libitum eligitur, ex $P$, $F$; ex $Q$, $G$; ex $R$, $H$ etc.: inter has formas $F,G,H$ etc. duae quae proprie aequivaleant esse non poterunt. Quaevis autem alia forma eiusdem determinantis alicui ex istis proprie aequivalens erit et quidem \textit{unicae} tantum. Hinc manifestum est, \textit{omnes formas huius determinantis in totidem classes distribui posse, quot habeantur periodi}, scilicet referendo eas quae formae $F$ proprie aequivalent in primam classem, eas quae formae $G$ proprie aequivalent in secundam etc. Hoc modo omnes formae in eadem classe contentae proprie aequivalentes erunt, formae vero e classibus diversis non poterunt proprie aequivalere. Sed hic huic argumento infra fusius explicando non immoramur.

\Needspace{8\baselineskip}
\Article{196}
\textsc{Problema.} \textit{Propositis duabus formis proprie aequivalentibus $\Phi$, $\varphi$: invenire transformationem propriam alterius in alteram.}

\textit{Sol.} Per methodum art. 183 inveniri poterunt duae series formarum
\[
\Phi,\Phi',\Phi''\mathinner{\ldotp\ldotp\ldotp}\Phi^n\quad\text{et}\quad\varphi,\varphi',\varphi''\mathinner{\ldotp\ldotp\ldotp}\varphi^\nu
\]
tales ut quaevis forma sequens praecedenti proprie aequivaleat, ultimaeque $\Phi^n$, $\varphi^\nu$ sint formae reductae; et quum $\Phi$, $\varphi$ proprie aequivalentes esse supponantur, necessario $\Phi^n$ in periodo formae $\varphi^\nu$ contenta erit. Sit $\varphi^\nu=f$ ipsiusque periodus usque ad formam $\Phi^n$ haec
\[
f,\ f',\ f''\mathinner{\ldotp\ldotp\ldotp} f^{m-1},\ \Phi^n
\]
ita ut in hac periodo index formae $\Phi^n$ sit $m$; designenturque formae quae oppositae sunt sociis formarum
\begin{center}
$\Phi,\ \Phi',\ \Phi''\mathinner{\ldotp\ldotp\ldotp}\Phi^n$\quad per $\Psi,\ \Psi',\ \Psi''\mathinner{\ldotp\ldotp\ldotp}\Psi^n$ resp.\HistoricalNote{*)}{Ita ut $\Psi$ oriatur ex $\Phi$ commutando terminum primum et ultimum tribuendoque medio signum oppositum, similiterque de reliquis.}.
\end{center}
Tum in progressione
\SourceLeaf{191}
\[
\varphi,\ \varphi',\ \varphi''\mathinner{\ldotp\ldotp\ldotp} f,\ f',\ f''\mathinner{\ldotp\ldotp\ldotp} f^{m-1},\ \Psi^{n-1},\ \Psi^{n-2}\mathinner{\ldotp\ldotp\ldotp}\Psi,\ \Phi
\]
quaevis forma praecedenti ab ultima parte contigua erit, unde per art. 177 inveniri poterit transformatio propria primae $\varphi$ in ultimam $\Phi$. Illud autem de formis reliquis progressionis nullo negotio perspicitur; de his $f^{m-1}$, $\Psi^{n-1}$ sic probatur: Sit
\[
\begin{gathered}f^{m-1}=(g,h,i);\quad f^m\text{ sive }\Phi^n=(g',h',i');\\ \Phi^{n-1}=(g'',h'',i'')\end{gathered}
\]
Forma $(g',h',i')$ tum formae $(g,h,i)$ tum formae $(g'',h'',i'')$ ab ultima parte contigua erit; hinc $i=g'=i''$, et $-h\equiv h'\equiv-h''\ (\mathrm{mod.}\ i\text{ vel }g'\text{ vel }i'')$. Unde manifestum est, formam $(i'',-h'',g'')$, \textit{i. e.} formam $\Psi^{n-1}$ formae $(g,h,i)$, \textit{i. e.} formae $f^{m-1}$ ab ultima parte contiguam esse.

Si formae $\Phi$, $\varphi$ improprie aequivalentes sunt: forma $\varphi$ proprie aequivalebit formae cui opposita est $\Phi$. Inveniri poterit itaque transformatio propria formae $\varphi$ in formam cui $\Phi$ est opposita; quae si supponitur fieri per substitutionem $\alpha,\beta,\gamma,\delta$, facile perspicitur, $\varphi$ improprie transformari in ipsam $\Phi$ per substitutionem $\alpha,-\beta,\gamma,-\delta$.

Hinc etiam perspicuum est, si formae $\Phi$, $\varphi$ tum proprie tum improprie aequivalentes sint, inveniri posse \textit{duas} transformationes, propriam et impropriam.

\textit{Ex.} Quaeritur transformatio impropria formae $(129,92,65)$ in formam $(42,59,81)$, quam illi improprie aequivalere in art. praec. invenimus. Investiganda erit itaque primo transformatio propria formae $(129,92,65)$ in formam $(42,-59,81)$. Ad hunc finem evolvitur progressio formarum haec: $(129,92,65)$, $(65,-27,10)$, $(10,7,-3)$, $(-3,8,5)$, $(5,22,81)$, $(81,59,42)$, $(42,-59,81)$. Hinc deducitur transformatio propria $-47,56,73,-87$, per quam $(129,92,65)$ transit in $(42,-59,81)$; quare per impropriam $-47,\allowbreak-56,\allowbreak73,\allowbreak87$ transibit in $(42,59,81)$.

\Needspace{8\baselineskip}
\Article{197}
Si transformatio una formae alicuius $(a,b,c)\mathinner{\ldotp\ldotp\ldotp}\varphi$ in aequivalentem $\Phi$ habetur: ex hac \textit{omnes} transformationes similes formae $\varphi$ in $\Phi$ deduci poterunt, si modo omnes solutiones aequationis indeterminatae $tt-Duu=mm$ assignari
